\scanpage{118}

\begin{flushright}
27/3/81
\end{flushright}

Cher Schurik,

Je n'ai pas de référence à proposer pour les pro-groupes libres, et l'essentiel de ce que je connais vient du Ch I de Serre, Cohomologie Galoisienne [cité CG]

On a

-- un sous-groupe d'un groupe libre est libre

\quad (clair par l'interprétation libre $\Longleftrightarrow\pi_1$ (un CW complexe de dim 1)

-- on comprend assez bien les pro-p-groupes libres, voire même les groupes nilpotents libres (?)

\margin{\figure{\begin{tikzpicture}[x=1cm,y=1cm]\draw[->] (0,.35) .. controls (.1,.15) and (.22,.05) .. (.25,0);\end{tikzpicture}}
voir CG, I5-I8 et références à Lazard}

et ces deux remarques doivent suffire pour tes questions sur ``centre'' et ``ordre fini''

a) si $u\in\hat L$ est d'ordre fini, soit $H\subset\hat L$ distingué ouvert ne contenant pas $u$\unsure{$\overset{\circ}{!}$}.
\margin{(mieux : $H\cap\mathbb Zu=\{e\}$)}
Remplaçant $L$ par $L\cap(H.\mathbb Zu)$, on se ramène à supposer \deleted{que} l'existence de
\[
L\overset{\varphi}{\twoheadrightarrow}\mathbb Z/n\mathbb Z
\quad\text{tel que }1\in\mathbb Z/n\mathbb Z\text{ puisse se relever}
\]
en un élément d'ordre $n$ \inserted{de $\hat L$}. Ceci est déjà faux avec $L$ remplacé par $L/(L,L)$.
\[
[(,)\text{ commutateurs}]
\]

b) pour $u\in\hat L$ central, on se ramène de même à supposer qu'on a
\[
L\overset{\varphi}{\twoheadrightarrow}\mathbb Z/n\mathbb Z,\quad
\varphi(u)\ne0,
\]
et une contradiction en regardant ce que serait l'image de \deleted{\unsure{$L$}} $u$ dans
\[
L/(\operatorname{Ker}\varphi,\operatorname{Ker}\varphi)
\qquad
\text{(et l'action de }u\text{ sur }
\operatorname{Ker}\varphi/(\operatorname{Ker}\varphi,\operatorname{Ker}\varphi))
\]

\margin{118}
